% ============================================================
%  09_validation_md.tex — Section 9 : Validation MD
% ============================================================
\section{Validation par dynamique moléculaire (MD)}

\dateblock{1\textsuperscript{er} mai 2026}

\begin{attention}
\textbf{Mise à jour du 6 mai 2026 — bascule vers un protocole déterministe.}
La phase MD elle-même a révélé un non-déterminisme inacceptable : xTB seede
les vitesses initiales sur l'horloge système (option \texttt{seed} dans le
bloc \texttt{\$md} ignorée), donc deux runs identiques produisaient des
trajectoires différentes — verdicts qui basculaient entre PLAN et NON PLAN
sur structures tendues (cas h6 alternés 5/7).

Décision arrêtée avec l'utilisateur : la priorité est la \textbf{reproductibilité
byte-pour-byte} (même \texttt{(graph, sizes)} $\to$ même \texttt{md\_final\_opt.xyz}),
quitte à perdre l'exploration thermique. La phase MD est remplacée par une
\textbf{perturbation z structurée déterministe} suivie de \texttt{xtb -{}-opt} :
\[
  z'_i = z_i + 0{,}05 \cdot \sin\!\left(\frac{2\pi i}{N} + 0{,}5\right)
\]
où $i$ est l'index canonique de l'atome (garanti par \texttt{sorted(vertex\_set)}
dans la reconstruction). Aucun RNG, aucune dépendance temporelle. Avec
\texttt{OMP\_NUM\_THREADS=1} + \texttt{MKL\_NUM\_THREADS=1}, md5 identique
sur N runs successifs. Vérifié sur 4 cas borderline h6 $\times$ 5 runs.

L'API publique \texttt{md\_then\_optimize(...)} et les noms de fichiers
de sortie (\texttt{md\_geom.xyz}, \texttt{md\_traj.xyz}, \texttt{md\_final\_opt.xyz})
sont conservés pour rétro-compat. Le bloc \texttt{md.inp} émis mentionne
désormais \texttt{method=det-opt} au lieu des anciens paramètres MD.

Le reste de cette section décrit le protocole MD historique
(avant 6 mai 2026), conservé pour traçabilité de la conception.
\end{attention}

\subsection{Motivation : protocole recommandé par les chimistes}

Après réunion avec Yannick Carissan (LS2N) et discussion sur les cas marginaux du
multi-runs, un \textbf{protocole alternatif} a été proposé :
\textit{lancer une dynamique moléculaire courte avant l'optimisation finale}.

\begin{quote}
\textit{Idée :} la reconstruction 3D produit une géométrie 2D plane par construction.
Une optimisation directe (\texttt{--opt}) reste piégée près de ce minimum local plat
même quand un minimum non-plat existe. La MD courte (1~ps à 298~K) injecte de
l'énergie thermique qui déforme la structure hors du plan, puis l'optimisation finale
relaxe vers le \emph{vrai} minimum énergétique. Si la structure y revient plate, elle est
réellement plate. Si elle se courbe, elle ne l'est pas.
\end{quote}

Ce protocole donne un verdict \textbf{binaire} (plan / non plan) sur un seul run par solution,
mais avec une validation physique plus stricte que le multi-runs perturbé.

\subsection{Protocole en 3 étapes}

Pour chaque solution issue du CSP :

\begin{enumerate}
  \item \texttt{xtb solution.xyz -{}-md} : MD NVT à 298~K, 1~ps, dump 50~fs $\to$ \texttt{md\_traj.xyz}
  \item Extraire la \textbf{dernière frame} de la trajectoire $\to$ \texttt{md\_geom.xyz}
  \item \texttt{xtb md\_geom.xyz -{}-opt tight} $\to$ \texttt{md\_final\_opt.xyz}, géométrie testée par l'ACP
\end{enumerate}

\subsection{Paramètres MD adoptés}

\begin{center}
\begin{tabular}{lll}
  \toprule
  Paramètre & Valeur & Justification \\
  \midrule
  Méthode quantique & GFN2-xTB    & cohérent avec le pipeline existant \\
  Ensemble          & NVT (Berendsen) & équilibrage thermique rapide \\
  Température       & 298.15~K    & condition standard \\
  Durée             & 1~ps        & suffisant pour s'écarter d'un plan métastable \\
  Pas               & 4~fs        & avec \texttt{hmass=4} pour relaxer la contrainte H \\
  \texttt{hmass}    & 4           & masse des H artificiellement augmentée \\
  Dump              & 50~fs       & $\sim$20 frames sur la trajectoire \\
  \texttt{velo}     & false       & xTB régénère les vélocités initiales \\
  \texttt{sccacc}   & 2.0         & tolérance SCC relâchée pour la MD (xTB) \\
  \bottomrule
\end{tabular}
\end{center}

\subsection{Pattern Strategy : décorrélation des méthodes de validation}

Pour permettre la \textbf{coexistence} des deux méthodes (multi-runs et MD) sans conflit, le
pipeline a été refactoré selon le \emph{Strategy Pattern}. Une classe abstraite
\texttt{ValidationStrategy} définit l'interface, deux implémentations concrètes la remplissent.

\begin{lstlisting}[style=benzai]
csp_solver/utils/validation/
  base.py              # ABC ValidationStrategy.validate_solutions(graph, sols, out_dir)
  multi_runs.py        # MultiRunsStrategy : delegue vers _run_single / _run_multi
  md.py                # MDStrategy : appelle optimizer_md.run_md_pipeline
  __init__.py          # Registry : get_strategy(name, **kwargs)
\end{lstlisting}

\begin{lstlisting}
class ValidationStrategy(ABC):
    name: str = ""
    @abstractmethod
    def validate_solutions(self, graph, solutions, output_dir):
        raise NotImplementedError

class MDStrategy(ValidationStrategy):
    name = "md"
    def __init__(self, threshold=10.0, opt_level="tight",
                 md_params=None, deterministic=True):
        ...
    def validate_solutions(self, graph, solutions, output_dir):
        # appelle non_benzenoid_generator/core/optimizer_md.py::run_md_pipeline
        ...
\end{lstlisting}

\texttt{main.py} accepte le drapeau \texttt{-{}-method} qui sélectionne la stratégie :
\texttt{md} ou \texttt{multi-runs}.

\dateblock{mai 2026 (mi-mois) -- defaut bascule sur \texttt{md}}

Depuis le bilan h3-h5 (98\% d'accord MR/MD) et la régénération de la table de voisinage
au protocole MD, \textbf{MD est désormais la méthode par défaut}. Le défaut historique
\texttt{multi-runs} reste accessible via \texttt{-{}-method multi-runs}, mais la
recommandation est d'utiliser MD pour tout nouveau run.

\begin{center}
\begin{tabular}{ll}
  \toprule
  Commande & Méthode utilisée \\
  \midrule
  \texttt{python main.py f.graph -{}-validate}                   & MD (défaut) \\
  \texttt{python main.py f.graph -{}-validate -{}-method md}     & MD (explicite) \\
  \texttt{python main.py f.graph -{}-validate -{}-method multi-runs -{}-n-runs 10} & multi-runs N=10 \\
  \bottomrule
\end{tabular}
\end{center}

\begin{attention}
  \texttt{-{}-n-runs N} est ignoré si \texttt{-{}-method multi-runs} n'est pas explicitement
  passé. C'est un changement de comportement par rapport à avril 2026 où \texttt{-{}-n-runs}
  suffisait pour basculer en multi-runs.
\end{attention}

\subsection{\texttt{aggregate\_md.py} — agrégation additive}

À l'image de \texttt{aggregate\_runs.py}, un script dédié recompose le \texttt{data.json}
d'une config à partir des fichiers \texttt{md\_validation/} produits par les runs MD :

\begin{lstlisting}[style=benzai]
output/h4/default/0-5-6-11/solutions/sol_1_5_7_7_5/
  md_validation/
    md_geom.xyz              # frame finale de la trajectoire (input de l'opt)
    md_traj.xyz              # trajectoire complete (input pour debug)
    md_final_opt.xyz         # geometrie testee par l'ACP
\end{lstlisting}

\textbf{Choix architectural important :} \texttt{aggregate\_md.py} est \textbf{additif}. Il
charge le \texttt{data.json} existant, ajoute un bloc \texttt{md\_validation} à chaque
solution sans toucher au bloc \texttt{runs} préexistant, puis ré-écrit. Cela permet d'avoir
sur le même dataset les deux méthodes côte à côte si on les a toutes deux lancées.

\dateblock{mai 2026 (mi-mois) -- reconciliation disque}

Le caractère additif a un effet de bord : si \texttt{data.json} contient des solutions
qui n'existent plus sur disque (par exemple parce qu'un re-run avec le nouveau défaut
exclut désormais la solution tout-hexagones, ou parce que la table régénérée laisse
moins de solutions), ces entrées resteraient indéfiniment dans le \texttt{data.json}.
Pour éviter ce drift, \texttt{aggregate\_md.py} \textbf{réconcilie} maintenant la
liste des solutions avec l'état du disque :

\begin{itemize}
  \item Pour chaque molécule, scan des dossiers \texttt{solutions/sol\_*/} pour
    récupérer les \texttt{sizes} effectivement présents (depuis le commentaire de
    \texttt{source.xyz}).
  \item Filtrage de \texttt{data["molecules"][mol]["solutions"]} : on ne garde que les
    entrées dont le \texttt{sizes} matche un dossier sur disque.
  \item Drop des molécules entières si leur dossier \texttt{mol/} n'existe plus.
\end{itemize}

Le même mécanisme symétrique est en place dans \texttt{aggregate\_runs.py} : avant
de réécrire \texttt{data.json}, on charge l'ancien et on \textbf{préserve les blocs
md\_validation} existants pour les solutions toujours présentes (matching par
\texttt{(mol\_name, sizes)}). Ainsi un re-run multi-runs sur une config qui avait déjà
des données MD ne perd pas ces données.

Schéma enrichi du \texttt{data.json} (extrait) :

\begin{lstlisting}[style=benzai]
{
  "molecules": {
    "0-5-6-11": {
      "solutions": [
        {
          "sizes": "v0=5 v1=7 v2=7 v3=5",
          "file": "sol_1_5_7_7_5/md_final_opt.xyz",
          "planar": true, "angle_deg": 0.4,
          "runs":          { ... },  /* multi-runs (si lance) */
          "md_validation": {         /* MD (si lance) */
            "planar": true,
            "angle_deg": 0.4,
            "rmsd": 0.03,
            "height": 0.12,
            "params": {"temp": 298.15, "time": 1.0, "step": 4, "dump": 50,
                       "hmass": 4, "sccacc": 2.0, "velo": false, "nvt": true},
            "trajectory_file": "md_validation/md_traj.xyz",
            "final_opt_file":  "md_validation/md_final_opt.xyz"
          }
        }
      ]
    }
  }
}
\end{lstlisting}

\subsection{Non-déterminisme contrôlé}

\texttt{xtb} initialise les vélocités MD à partir de l'horloge système : deux exécutions
identiques produisent des trajectoires différentes. Ce comportement n'est pas configurable
côté utilisateur (l'option \texttt{seed} dans le bloc \texttt{\$md} est ignorée par xTB).

\textbf{Mitigation} : \texttt{MDStrategy} fixe \texttt{OMP\_NUM\_THREADS=1} et
\texttt{MKL\_NUM\_THREADS=1} dans l'environnement du subprocess. Cela élimine
le non-déterminisme de l'addition flottante en parallèle, ce qui réduit la variance
des angles finaux à environ \textbf{0.01°} (très en dessous du seuil 10°). Les verdicts
plan/non plan sont donc stables d'une exécution à l'autre, même si les trajectoires
diffèrent en détail.

Désactivable via \texttt{-{}-md-no-deterministic} si on veut récupérer le parallélisme.

\subsection{Lancement des batchs MD}

Le drapeau \texttt{-{}-method md} se propage de \texttt{batch\_all.py} jusqu'à
\texttt{main.py}, et déclenche \texttt{aggregate\_md.py} en fin de batch (au lieu
d'\texttt{aggregate\_runs.py}) :

\begin{lstlisting}[style=benzai]
# Batch MD complet (8 configs) sur h3
python batch_all.py plane/benzdb/h3 --method md

# Une seule config MD
python batch_main.py plane/benzdb/h3 --validate --method md --no-freeze --adj-57
\end{lstlisting}

\subsection{Comparaison MR vs MD : résultats}

Sur le dataset h3--h5 (53 solutions, 8 configs), les deux méthodes convergent à
\textbf{98\%} : 52 verdicts plan/non-plan identiques, 1 seul cas de divergence franche
(MR plan, MD non plan ou inverse) sur la config \texttt{default}.

Les multi-runs apportent une \emph{nuance qualitative} (always\_/mostly\_planar,
unstable, etc.) mais identifient le même ensemble de cas problématiques que le MD.
Le MD est plus court à exécuter (1 run au lieu de 10) mais physiquement plus rigoureux.
